\documentclass[a4paper,10pt]{article}
\usepackage[utf8]{inputenc}
\usepackage{geometry}
\geometry{left=1in, right=1in, top=1in, bottom=1in}

\title{Revised Prioritized Project and Study Plan}
\author{Stephen MacKenzie}
\date{}

\begin{document}

\maketitle

\section*{Prioritized Plan of Action}

\subsection*{1. Highest Priority: MLibs Project Video}
\begin{itemize}
    \item \textbf{Objective}: Create and upload a video demonstrating the MLibs project, including its setup, components, and real-time execution.
    \item \textbf{Tasks}:
    \begin{itemize}
        \item Prepare a video script outlining the introduction, project overview, and real-time execution of the MLibs project.
        \item Record the setup, build, and execution process using screen recording software.
        \item Include an explanation of each library module (priority queues, hash tables, etc.) with Doxygen documentation as visual aids.
        \item Demonstrate the automated build, flash, and run process using the \texttt{bfr.py} script.
        \item Upload the completed video to YouTube and share it on LinkedIn.
    \end{itemize}
    \item \textbf{Timeframe}: 1-2 days
\end{itemize}

\subsection*{2. Real-Time Packet Reading and Interaction (Voltage Analyzer Project)}
\begin{itemize}
    \item \textbf{Objective}: Enhance the \texttt{voltage\_analyzer} project to read Ethernet packets in real-time from the microcontroller and interact with the microcontroller using two-way communication.
    \item \textbf{Tasks}:
    \begin{itemize}
        \item Implement a thread to handle real-time packet reading using non-blocking socket I/O and `ncurses` for display.
        \item Set up a thread for sending packets to the microcontroller, triggering interrupts for polling data.
        \item Implement interrupt handling and threading to synchronize data processing.
        \item Integrate temperature monitoring code, making the application a comprehensive remote monitor and analyzer.
        \item Add Doxygen support to generate documentation for each module and provide detailed comments for better code understanding.
    \end{itemize}
    \item \textbf{Outcome}: A real-time remote monitor and analyzer that reads and displays voltage and temperature data from the microcontroller, showcasing multi-threaded interactions and low-level network programming.
    \item \textbf{Timeframe}: 1 week
\end{itemize}

\subsection*{3. Multithreaded Development Assignment}
\begin{itemize}
    \item \textbf{Objective}: Implement a multi-threaded data pipeline that utilizes various threading primitives for synchronization and coordination.
    \item \textbf{Tasks}:
    \begin{itemize}
        \item Develop a data processing pipeline with multiple threads, each handling different stages of data transformation (e.g., reading, processing, and aggregating).
        \item Implement synchronization using \texttt{std::mutex}, \texttt{std::atomic}, and \texttt{std::condition\_variable}.
        \item Add demonstrations of thread-safe data structures and synchronization methods.
        \item Integrate \texttt{std::numerics} components for computational tasks like FFT.
    \end{itemize}
    \item \textbf{Outcome}: A comprehensive multi-threaded project that showcases your understanding of synchronization and parallelism in C++.
    \item \textbf{Timeframe}: 2-3 weeks
\end{itemize}

\subsection*{4. Numerics Library Integration}
\begin{itemize}
    \item \textbf{Objective}: Understand and implement numerics operations using \texttt{std::numeric} and related libraries to complement threading work.
    \item \textbf{Tasks}:
    \begin{itemize}
        \item Explore and implement key numerics operations (e.g., \texttt{std::valarray}, \texttt{std::complex}, \texttt{std::accumulate}).
        \item Integrate numerics computations into the threading project (e.g., using FFT for data processing).
        \item Profile and optimize numerics operations in multi-threaded scenarios.
    \end{itemize}
    \item \textbf{Outcome}: A strong grasp of numerics utilities and how they can be used effectively in concurrent environments.
    \item \textbf{Timeframe}: 1-2 weeks
\end{itemize}

\subsection*{5. CLRS-Based Assignments (Lower Priority)}
\begin{itemize}
    \item \textbf{Objective}: Reinforce algorithmic problem-solving skills with targeted exercises from CLRS.
    \item \textbf{Tasks}:
    \begin{itemize}
        \item Implement dynamic programming exercises like the Rod Cutting problem.
        \item Work on graph theory problems such as BFS, DFS, and concurrent graph traversal algorithms.
        \item Implement concurrent versions of algorithms (e.g., parallel Dijkstra's shortest path) using threading concepts.
    \end{itemize}
    \item \textbf{Outcome}: A solid understanding of algorithms that complement your threading and numerics work.
    \item \textbf{Timeframe}: As time permits, revisit after completing the primary threading and numerics work.
\end{itemize}

\subsection*{6. Additional Stretch Goals (Optional)}
\begin{itemize}
    \item \textbf{Objective}: Implement advanced projects that leverage custom allocators, memory management, and real-world simulations.
    \item \textbf{Tasks}:
    \begin{itemize}
        \item Implement a memory pool allocator for multi-threaded allocation and deallocation.
        \item Simulate real-world scenarios like traffic management using threading and graph algorithms.
    \end{itemize}
    \item \textbf{Outcome}: A deeper understanding of memory management and concurrency in real-world systems.
    \item \textbf{Timeframe}: Optional, if time permits.
\end{itemize}

\end{document}
